% !TeX program = pdflatex
\documentclass[11pt,a4paper,oneside]{article}
\usepackage[T1]{fontenc}
\usepackage[utf8]{inputenc}
\usepackage[main=french,provide=*]{babel}
\usepackage{lmodern,microtype}
\usepackage{amsmath,amssymb,mathtools}
\usepackage{geometry}
\geometry{left=2.3cm,right=2.3cm,top=2.55cm,bottom=2.35cm,headheight=24pt}
\usepackage{array,booktabs,tabularx}
\usepackage{enumitem}
\usepackage{xcolor}
\usepackage{tikz}
\usetikzlibrary{calc}
\usepackage[most]{tcolorbox}
\usepackage{fancyhdr,lastpage}
\usepackage{titlesec}
\usepackage{needspace}
\usepackage{hyperref}
\definecolor{Marine}{HTML}{10283F}
\definecolor{Bleu}{HTML}{176D8E}
\definecolor{Sarcelle}{HTML}{118F86}
\definecolor{Or}{HTML}{EDAA47}
\definecolor{GrisTexte}{HTML}{54667A}
\definecolor{BleuPale}{HTML}{EFF5F8}
\definecolor{OrangePale}{HTML}{FFF6E7}
\hypersetup{colorlinks=true,linkcolor=Bleu,urlcolor=Sarcelle,
pdftitle={Algèbre L1 Chimie - Chapitre 2},
pdfauthor={M. Bakary TRAORE - Universite de N'Zerekore}}
\setlength{\parindent}{0pt}
\setlength{\parskip}{4pt plus1pt minus1pt}
\setlist[enumerate]{leftmargin=1.6em,itemsep=2pt,topsep=3pt}
\setlist[itemize]{leftmargin=1.4em,itemsep=2pt,topsep=3pt}
\renewcommand{\arraystretch}{1.15}
\emergencystretch=2em
\titleformat{\section}{\Large\bfseries\color{Marine}}{\thesection}{.65em}{}[\vspace{1mm}{\color{Sarcelle}\titlerule[.75pt]}]
\titleformat{\subsection}{\normalsize\bfseries\color{Bleu}}{}{0pt}{}
\titlespacing*{\section}{0pt}{18pt}{9pt}
\titlespacing*{\subsection}{0pt}{10pt}{4pt}
\pagestyle{fancy}
\fancyhf{}
\fancyhead[L]{\small\color{Marine}\textbf{UNIVERSITÉ DE N'ZÉRÉKORÉ}}
\fancyhead[R]{\small\color{GrisTexte}Algèbre \textbullet\ 01-MTH 1102}
\fancyfoot[L]{\scriptsize\color{GrisTexte}Chapitre 2 \textbullet\ L1 Chimie}
\fancyfoot[R]{\scriptsize\color{GrisTexte}\thepage\,/\,\pageref{LastPage}}
\renewcommand{\headrulewidth}{.5pt}
\renewcommand{\footrulewidth}{.3pt}
\renewcommand{\headrule}{\hbox to\headwidth{\color{Sarcelle}\leaders\hrule height \headrulewidth\hfill}}
\renewcommand{\footrule}{\hbox to\headwidth{\color{Bleu!30}\leaders\hrule height \footrulewidth\hfill}}
\newtcolorbox{objectifbox}[1][]{enhanced,breakable,colback=BleuPale,colframe=Bleu,
boxrule=.6pt,arc=1mm,#1}
\newtcolorbox{conseilbox}[1][]{enhanced,breakable,colback=OrangePale,colframe=Or!75!black,
boxrule=.6pt,arc=1mm,#1}
\newcommand{\exo}[3]{%
  \Needspace{6\baselineskip}%
  \subsection*{Exercice #1 --- #3 \hfill
  {\normalfont\footnotesize\color{GrisTexte}[#2]}}}
\newcommand{\R}{\mathbb{R}}
\newcommand{\D}{\mathcal{D}}

\newcommand{\C}{\mathbb C}
\newcommand{\ii}{\mathrm i}
\newcommand{\ee}{\mathrm e}
\newtcolorbox{defbox}[1]{enhanced,breakable,colback=BleuPale,colframe=Bleu,boxrule=.6pt,arc=1mm,title={#1},fonttitle=\bfseries}
\newtcolorbox{exbox}[1]{enhanced,breakable,colback=white,colframe=Sarcelle,boxrule=.6pt,arc=1mm,title={#1},fonttitle=\bfseries}
\begin{document}
\hypersetup{pageanchor=false}
\thispagestyle{empty}
\begin{tikzpicture}[remember picture,overlay]
\fill[Bleu] (current page.south west) rectangle (current page.north east);
\draw[white,opacity=.09,line width=1pt] ([xshift=2.6cm,yshift=5cm]current page.south west) circle (4.1cm);
\draw[white,opacity=.09,line width=1pt] ([xshift=16.2cm,yshift=23.5cm]current page.south west) circle (5.4cm);
\draw[Or,opacity=.8,line width=2pt] ([xshift=3cm,yshift=7.2cm]current page.south west) -- ([xshift=18cm,yshift=7.2cm]current page.south west);
\end{tikzpicture}
\color{white}\vspace*{2.85cm}
\begin{center}
{\fontsize{19}{23}\selectfont\bfseries UNIVERSITÉ DE N'ZÉRÉKORÉ\par}
\vspace{.5cm}{\large LICENCE 1 EN CHIMIE\par}
\vspace{2.1cm}
{\color{Or}\fontsize{20}{25}\selectfont\bfseries TRAVAUX DIRIGÉS 2\par}
\vspace{.75cm}{\fontsize{34}{40}\selectfont\bfseries ALGÈBRE\par}
\vspace{.45cm}{\fontsize{25}{31}\selectfont\bfseries CHAPITRE 2\par}
\vspace{.5cm}{\fontsize{20}{25}\selectfont Nombres complexes et polynômes\par}
\vspace{.5cm}{\large Complexes \textbullet\ Polynômes \textbullet\ Chimie\par}
\vspace{4.5cm}{\normalsize Enseignant : \textbf{M. Bakary TRAORE}\par}
\vspace{.35cm}{\normalsize Année universitaire 2026--2027\par}
\end{center}
\clearpage\color{black}\pagenumbering{arabic}
\hypersetup{pageanchor=true}\setcounter{page}{1}


\section*{Consignes et organisation}
\begin{objectifbox}
\textbf{Durée :} 2 h de TD. Les exercices 1 à 6 constituent le parcours principal (environ 100 minutes) ; les exercices 7 et 8 prolongent l'entraînement (20 minutes en classe ou à domicile). Justifier chaque réponse, préciser le corps de coefficients, et vérifier les calculs par substitution ou développement. Les angles sont en radians.
\end{objectifbox}
\section{Nombres complexes}
\exo{1}{Fondamental}{Forme algébrique et opérations}
Soient $z=3-2\ii$ et $w=-1+4\ii$.
\begin{enumerate}[label=\alph*)]
\item Donner les parties réelle et imaginaire de $z$ et $w$.
\item Calculer $z+w$, $z-w$ et $zw$ sous forme algébrique.
\item Calculer $\ii^{37}$ et $(1-\ii)^2$.
\end{enumerate}
\exo{2}{Fondamental}{Conjugué, module et quotient}
Soient $u=2+3\ii$ et $v=1-\ii$.
\begin{enumerate}[label=\alph*)]
\item Calculer $\overline u$, $|u|$ et $|v|$ ; vérifier $u\overline u=|u|^2$.
\item Calculer $u/v$ sous forme algébrique.
\item Vérifier votre quotient par multiplication, puis calculer son module de deux façons.
\end{enumerate}
\exo{3}{Intermédiaire}{Plan complexe et formes polaires}
Soient $z=-\sqrt3+\ii$ et $w=1+\ii$.
\begin{enumerate}[label=\alph*)]
\item Placer leurs points dans le plan complexe et déterminer les modules et un argument principal de chacun.
\item Écrire les formes trigonométrique et exponentielle de $z$ et $w$.
\item Donner les formes exponentielles de $zw$ et $z/w$. Calculer $w^8$ sous forme algébrique.
\end{enumerate}
\section{Polynômes}
\exo{4}{Fondamental}{Calcul polynomial}
Soient $P=2X^2-X+3$ et $Q=X^2+2X-1$.
\begin{enumerate}[label=\alph*)]
\item Donner le degré et le coefficient dominant de chacun.
\item Calculer $P+Q$, $P-2Q$ et $PQ$ ; préciser leurs degrés.
\item Calculer $P(-1)$ puis évaluer $H(X)=2X^3-X^2+3X-4$ en $X=2$ par Horner.
\end{enumerate}
\exo{5}{Intermédiaire}{Division euclidienne et reste}
\begin{enumerate}[label=\alph*)]
\item Effectuer la division euclidienne de $A=X^3+3X^2+2X+5$ par $B=X^2+2$.
\item Vérifier l'identité $A=BQ+R$ et la condition sur le degré du reste.
\item Sans effectuer toute la division, déterminer le reste de la division de $A$ par $X-1$. Ce dernier divise-t-il $A$ ?
\end{enumerate}
\exo{6}{Intermédiaire}{Racines et factorisation}
\begin{enumerate}[label=\alph*)]
\item Résoudre $X^2+2X+5=0$ dans $\R$ puis dans $\C$. Factoriser dans $\C[X]$.
\item Soit $F=X^3-5X^2+8X-4$. Montrer que $2$ est une racine, puis factoriser $F$ complètement dans $\R[X]$.
\item Donner les multiplicités des racines de $F$ et vérifier en développant le produit.
\end{enumerate}
\section{Applications et approfondissement}
\exo{7}{Application}{Calibration d'un instrument}
Un modèle fictif de calibration donne $S(c)=0{,}05c^2+0{,}3c$ pour $0\le c\le6$, où $c$ est exprimée en mmol/L et $S$ est sans unité.
\begin{enumerate}[label=\alph*)]
\item Une réponse $S=1{,}35$ est observée. Écrire une équation polynomiale équivalente à cette mesure.
\item Déterminer les racines ; choisir celle qui est compatible avec la plage de calibration et vérifier.
\item Une autre réponse vaut $4$. Est-elle compatible avec la plage annoncée ? Justifier sans dérivation, en encadrant $S(c)$.
\end{enumerate}
\exo{8}{Application}{Signal oscillant et racines complexes}
On représente un signal idéal par $Z(t)=5\ee^{\ii(4t+\pi/6)}$, avec $t$ en secondes.
\begin{enumerate}[label=\alph*)]
\item Exprimer le signal réel $s(t)=\Re(Z(t))$, donner son amplitude, sa phase initiale et sa période.
\item Calculer $Z(0)$ sous forme algébrique et $s(0)$.
\item En utilisant la forme exponentielle, résoudre $u^2=\ii$ dans $\C$ ; vérifier les deux solutions.
\end{enumerate}
\begin{conseilbox}Les valeurs numériques des applications sont fictives. Les exercices 7 et 8 illustrent les outils mathématiques et ne constituent pas des protocoles de laboratoire.
\end{conseilbox}

\end{document}
